% ==============================================================================
% annexes/about_author.tex — UE SIS589
% À Propos de l'Auteur
% ==============================================================================

\chapter*{À Propos de l'Auteur}
\addcontentsline{toc}{chapter}{À Propos de l'Auteur}
\label{ann:auteur}

\begin{center}
{\large\bfseries\color{CouleurPrimaire}
MINKA MI NGUIDJOÏ Thierry Emmanuel}\\[0.5ex]
{\normalsize\color{CouleurBordInfo}
Matricule~17P318~·~Chercheur LIMSI~·~Enseignant ENSP~·~Auditeur Senior CONSUPE}
\end{center}

\bigskip

% ==============================================================================
\section*{Profil académique et de recherche}
% ==============================================================================

MINKA MI NGUIDJOÏ Thierry Emmanuel est doctorant au Laboratoire
d'Informatique et Mathématiques pour les Systèmes Informatiques (LIMSI)
de l'École Nationale Supérieure Polytechnique (ENSP) de l'Université de
Yaoundé~I, sous la direction du Professeur Thomas Djotio Ndié et du
Docteur Flavien Serge Mani Onana. Sa thèse, intitulée \textit{«~Aux
Fondements de la Cryptographie Institutionnelle~»}, est en attente de soutenance. 
Elle apporte une contribution fondamentale à la formalisation
des systèmes cryptographiques en contexte d'audit institutionnel et d'État
de droit numérique.

Ses travaux de recherche publiés dans la littérature scientifique internationale
comprennent notamment~:
\begin{itemize}
  \item \textbf{P1~---} \textit{The CLO Framework~: Cryptographic Legitimacy Operators
        for Institutional Proof Systems}, IEEE Access, DOI~:
        10.1109/ACCESS.2026.3660105, pp.~20427--20445.
  \item \textbf{P2~---} \textit{The CRO Trilemma~: Completeness, Reliability and
        Opacity in Institutional Cryptography}, IEEE Access, DOI~:
        10.1109/ACCESS.2026.3669065, pp.~33263--33274.
  \item \textbf{S1~---} \textit{Post-Quantum CRO Extension} (sous révision).
  \item \textbf{S2~---} \textit{The SOTP Family~: Scalable One-Time Proof Systems}
        (sous révision).
  \item \textbf{S3~---} \textit{Formalizing the Conditional CRO Trilemma with Lean4}
        (sous révision).

  \item \textbf{p1~---} \textit{ZK-NR: A Layered Cryptographic Architecture for Explainable Non-Repudiation}, IACR eprint, https://eprint.iacr.org/2025/1138.
  \item \textbf{p2~---} \textit{The CRO Trilemma : a formal incompatibility between Confidentiality, Reliability and legal Opposability in Post-Quantum proof systems}, IACR eprint,
        https://eprint.iacr.org/2025/1348.
  \item \textbf{p3~---} \textit{Quantum Composable and Contextual Security Infrastructure (Q2CSI) : A Modular Architecture for Legally Explainable Cryptographic Signatures}, IACR eprint, https://eprint.iacr.org/2025/1380.
  \item \textbf{p4~---} \textit{Design ZK-NR: A Post-Quantum Layered Protocol for Legally Explainable Zero-Knowledge Non-Repudiation Attestation}, IACR eprint, https://eprint.iacr.org/2025/1422.
        (sous révision).
  \item \textbf{p5~---} \textit{UC-Security of the ZK-NR Protocol under Contextual Entropy Constraints: A Composable Zero-Knowledge Attestation Framework}, IACR eprint, https://eprint.iacr.org/2025/1529.

  \item \textbf{p6~---} \textit{The Chaotic Entropic Expansion (CEE): A Transparent Post-Quantum Data Confidentiality Primitive via Entropic Chaotic Maps}, IACR eprint, https://eprint.iacr.org/2025/1615.
   
  \item \textbf{p7~---} \textit{The AIIP Problem: Toward a Post-Quantum Hardness Assumption from Affine Iterated Inversion over Finite Fields}, IACR eprint, https://eprint.iacr.org/2025/1590.
  
  \item \textbf{p8~---} \textit{The Affine One-Wayness (AOW): A Transparent Post-Quantum Temporal Verification via Polynomial Iteration}, IACR eprint, https://eprint.iacr.org/2025/1662.
  \item \textbf{p9~---} \textit{The Semantic Holder (SH): Algebraic Extraction for Legal Opposability}, IACR eprint, https://eprint.iacr.org/2025/1708.
  
\end{itemize}


Il est également Research Director en Cybersécurité à
\textbf{AfroLeadership} (plateforme panafricaine de développement du
leadership), où il coordonne les initiatives de formation et de
sensibilisation à la sécurité numérique à l'échelle continentale.

% ==============================================================================
\section*{Certifications professionnelles}
% ==============================================================================

Thierry MINKA est titulaire de l'ensemble des certifications majeures
de l'\textbf{ISACA} (Information Systems Audit and Control Association),
confirmant son expertise à l'intersection de l'audit, de la gouvernance
et de la sécurité des systèmes d'information~:

\begin{table}[H]
\centering\small
\rowcolors{2}{TableauPair}{TableauImpair}
\begin{tabular}{p{2cm}p{8cm}p{3.5cm}}
\toprule
\tableauHeader{Sigle} & \tableauHeader{Intitulé complet} &
\tableauHeader{Organisme} \\
\midrule
\textbf{CISA} & Certified Information Systems Auditor & ISACA \\
\textbf{CISM} & Certified Information Security Manager & ISACA \\
\textbf{CGEIT} & Certified in the Governance of Enterprise IT & ISACA \\
\textbf{CRISC} & Certified in Risk and Information Systems Control & ISACA \\
\textbf{CDPSE} & Certified Data Privacy Solutions Engineer & ISACA \\
ISO~27001 & Lead Implementer / Lead Auditor & Bureau Veritas / PECB \\
ISO~22301 & Business Continuity Lead Auditor & PECB \\
ISO~9001 & Quality Management Lead Auditor & PECB \\
\bottomrule
\end{tabular}
\end{table}

% ==============================================================================
\section*{Expérience professionnelle}
% ==============================================================================

\textbf{Auditeur Senior — CONSUPE (Contrôle Supérieur de l'État du Cameroun).}
Thierry MINKA exerce en qualité d'auditeur senior au sein de la plus haute
institution d'audit interne de l'État camerounais. Il a conduit ou supervisé
des dizaines de missions d'audit des systèmes d'information dans les
administrations publiques et les entreprises à participation d'État.
Une de ses missions de terrain phares (2018--2019) a ciblé le système
\textbf{ANTILOPE} (plateforme de gestion de la solde des fonctionnaires
de l'État camerounais). Cette mission a mis en lumière des anomalies
cryptographiques fondamentales~: comment produire une preuve admissible
en justice sans compromettre le secret d'État~? Ce dilemme pratique
est à l'origine de toute la recherche en cryptographie institutionnelle
développée dans sa thèse.

\textbf{Stream Leader Cybersécurité \& Résilience — CICA-RE / Forvis Mazars.}
Dans le cadre de la modernisation du Système d'Information de la CICA-RE
(Compagnie Commune de Réassurance des États Membres de la CICA), il dirige
le Stream~3 portant sur la cybersécurité et la résilience, sous la supervision
du cabinet Forvis Mazars.

\textbf{Expert Judiciaire en Forensique Numérique.}
Inscrit sur la liste des experts judiciaires, il intervient en qualité
d'expert auprès des juridictions pour l'analyse de preuves numériques dans
des affaires de cybercriminalité, de fraude informatique et de litiges
portant sur des systèmes d'information.

% ==============================================================================
\section*{Enseignement}
% ==============================================================================

Thierry MINKA enseigne dans le supérieur comme enseignant vacataire depui 2016, notamment~:
À l'Université de Yaoundé I. (ENSP- Département de Génie Informatique), 
\begin{itemize}
  \item \textbf{UE INF3226~--} Audit des Systèmes d'Information.
  \item \textbf{UE SEC3266~--} Sécurité Informatique.
  \item \textbf{UE GES4062~--} Analyse de Risque et Gestion de Crise.
  \item \textbf{UE SEC4062~--} Hacking Éthique et Tests d'Intrusion.
  \item \textbf{UE SEC4031~--} Investigation Numérique I.
  \item \textbf{UE SEC4052~--} Investigation Numérique II.
  
\end{itemize}
À l'Université de Garoua (Département de Mathématiques), 
\begin{itemize}
  \item \textbf{UE LPI342~--} Forensique Numérique.
  \item \textbf{UE SEC3266~--} Systèmes de Gestion de la Sécurité Informatique.
  \item \textbf{UE SIS589~--} Supervision, Analyse Forensique et Gestion
        des Incidents. Le présent support de cours.
\end{itemize}

À l'Université de Yaoundé II (Programme Supérieur de Spécialisation en Finances Publiques), 
\begin{itemize}
  \item  Audit de Systèmes d'Information.

\end{itemize}

Sa pédagogie repose sur le fil rouge \textbf{LiZbiZ~/ LiZbiZ-SIS}~:
une entreprise fictive camerounaise dont l'infrastructure, les incidents
et les vulnérabilités évoluent de cours en cours, permettant aux étudiants
de construire progressivement une vision intégrée de la sécurité des SI,
de l'audit à la réponse aux incidents.

Il conduit également des \textbf{sessions de sensibilisation OSINT} où il
utilise sa propre empreinte numérique comme cas d'étude pédagogique,
démontrant concrètement les informations collectables sur un individu
par des techniques légales de sources ouvertes.

l'ensemble des ses cours est disponible sur la plateforme GitHub à l'adresse https://github.com/MaletYon.
% ==============================================================================
\section*{Contact}
% ==============================================================================

\begin{center}
\begin{minipage}{0.7\textwidth}
\centering
{\color{CouleurPrimaire}\bfseries MINKA MI NGUIDJOÏ Thierry Emmanuel}\\[0.5ex]
https://github.com/MaletYon.\\
maletyon@proton.me\\
\end{center}

\vspace{2ex}
\begin{center}
{\color{CouleurAccent}\rule{0.5\linewidth}{0.8pt}}\\[0.5ex]
{\small\itshape\color{CouleurBordInfo}
\og{}La cryptographie sans contexte institutionnel est un outil sans finalité.\\
La sécurité sans culture organisationnelle est une illusion sans durée.\fg{}}
\end{center}
